\section{Column dimensions and finite-group fixed spaces}
\label{cd-section}

We prove the dimension formula for the actual decorated source
space, including its finite-level factors and its coefficient
dimension. We also compute the fixed spaces and simultaneous
kernels of the norm operators at the auxiliary primes.
These are finite-dimensional algebraic statements. The mass which
appears below is defined as a weighted sum of lattice classes;
its separate analytic product formula is not used in this section.
The definitions and prime conditions are those of the submitted
column construction.\cite{alpo-uniform-source}

\subsection{Evaluation of columns on double cosets}

Let $\Gamma,G_\infty,G_f$ be groups, with homomorphisms
$\iota_\infty:\Gamma\to G_\infty$ and
$\iota_f:\Gamma\to G_f$. Let $K\subseteq G_f$ be a subgroup,
$k$ a field, $\psi:K\to k^\times$ a character, and
$\pi:G_\infty\to\operatorname{GL}(V)$ a representation on a
finite-dimensional $k$-space. A column is a linear map
$T:V\to\{\text{functions }G_\infty\times G_f\to k\}$ such that
\begin{align}
 (Tv)(\iota_\infty(\gamma)g,\iota_f(\gamma)x)
       &=(Tv)(g,x),\notag\\
 (Tv)(g,x\kappa)&=\psi(\kappa)(Tv)(g,x),\notag\\
 (T\pi(h)v)(g,x)&=(Tv)(gh,x)
 \label{cd-column-equivariance}
\end{align}
for all indicated group elements. Denote this vector space by
$\mathcal C(\pi,K,\psi)$.

\begin{proposition}\label{cd-column-decomposition}
Suppose the double-coset set
$\mathcal X=\iota_f(\Gamma)\backslash G_f/K$ is finite and
\begin{equation}\label{cd-trivial-stabilizers}
     x^{-1}\iota_f(\gamma)x\in K\quad\Longrightarrow\quad\gamma=1
                            \qquad(x\in G_f,\ \gamma\in\Gamma).
\end{equation}
Choose representatives $s_c$, $c\in\mathcal X$. Evaluation gives
a linear isomorphism
\[
 \mathcal C(\pi,K,\psi)\longrightarrow\bigoplus_{c\in\mathcal X}V^*,
 \qquad T\longmapsto\bigl(v\mapsto(Tv)(1,s_c)\bigr)_c.
\]
In particular
\begin{equation}\label{cd-column-rank}
                  \dim_k\mathcal C(\pi,K,\psi)
                               =|\mathcal X|\dim_k V.
\end{equation}
The same conclusion holds if $K$ is contained in a larger subgroup
for which \eqref{cd-trivial-stabilizers} is known.
\end{proposition}
\begin{proof}
Every $x\in G_f$ can be written
$x=\iota_f(\gamma)s_c\kappa$. This expression is unique once
$c$ and its representative are fixed. Indeed, two such expressions
give
\[
 s_c^{-1}\iota_f(\gamma'^{-1}\gamma)s_c=\kappa'\kappa^{-1}\in K,
\]
so \eqref{cd-trivial-stabilizers} gives $\gamma=\gamma'$ and
then $\kappa=\kappa'$.
For any collection $L_c\in V^*$ define
\begin{equation}\label{cd-column-extension}
       (Tv)(g,\iota_f(\gamma)s_c\kappa)
          =\psi(\kappa)L_c\bigl(\pi(\iota_\infty(\gamma)^{-1}g)v\bigr).
\end{equation}
Uniqueness makes this well-defined. The homomorphism laws show
directly that it satisfies all three identities
\eqref{cd-column-equivariance}. Those identities also force
\eqref{cd-column-extension} for any given column. Thus extension
and evaluation are inverse. Containment in a larger subgroup
preserves \eqref{cd-trivial-stabilizers}.
\end{proof}

There is a useful convention check over $\mathbb C$. Suppose $V$
has a Hermitian form linear in its first argument, $\pi$ is
unitary, and $|\psi|=1$. The unique vectors $F(x)$ satisfying
$(Tv)(1,x)=\langle v,F(x)\rangle$ obey
\begin{equation}\label{cd-riesz-character}
 F(\iota_f(\gamma)x\kappa)
            =\psi(\kappa)^{-1}\pi(\iota_\infty(\gamma))F(x).
\end{equation}
Indeed \eqref{cd-column-equivariance} gives
\[
 (Tv)(1,\iota_f(\gamma)x\kappa)
     =\psi(\kappa)\langle\pi(\iota_\infty(\gamma)^{-1})v,F(x)\rangle,
\]
and unitarity and conjugate linearity in the second argument give
\eqref{cd-riesz-character}. The passage $T\mapsto F$ is
conjugate linear, not complex linear. With the opposite Hermitian
convention, the same identity is written
$(Tv)(1,x)=\langle F(x),v\rangle$. This is the source's convention.

Without \eqref{cd-trivial-stabilizers}, the evaluated covectors
must satisfy the actual stabilizer conditions
\[
 L_c\bigl(\pi(\iota_\infty(\gamma))v\bigr)
   =\psi(s_c^{-1}\iota_f(\gamma)s_c)^{-1}L_c(v)
 \quad\text{if }s_c^{-1}\iota_f(\gamma)s_c\in K.
\]
They cannot in general be replaced by all of $V^*$.

\subsection{The fixed-space count with its character hypotheses}

Let now $\Gamma,K,K_0,D$ be subgroups of a group $G$, with
$K,D\subseteq K_0$. Assume
\begin{equation}\label{cd-neat}
                         \Gamma\cap xK_0x^{-1}=\{1\}
                                             \qquad(x\in G).
\end{equation}
Let $W$ be a finite-dimensional $k$-representation $\rho$ of
$\Gamma$, and let $\chi:K\to k^\times$ be a character.
Write $\mathcal A_K$ for the functions $f:G\to W$ satisfying
\[
                         f(\gamma x\kappa)
                                =\chi(\kappa)\rho(\gamma)f(x).
\]
Suppose $D$ normalizes $K$, $\chi(\delta\kappa\delta^{-1})
=\chi(\kappa)$ for $\delta\in D$, and $\chi$ is trivial on
$D\cap K$. Denote by $\mathcal A_K^D$ those functions which
also satisfy $f(x\delta)=f(x)$ for all $\delta\in D$.

\begin{proposition}\label{cd-fixed-dimension}
If $\Gamma\backslash G/K$ is finite, then $[D:D\cap K]$ is
finite and
\begin{equation}\label{cd-fixed-dimension-eq}
                   \dim_k\mathcal A_K^D\,[D:D\cap K]
                                     =\dim_k\mathcal A_K.
\end{equation}
This holds in every characteristic, without dividing by the index.
\end{proposition}
\begin{proof}
The product $L=DK=KD$ is a subgroup of $K_0$. Define
$\widetilde\chi(\delta\kappa)=\chi(\kappa)$.
If $\delta\kappa=\delta'\kappa'$, then
$\kappa'\kappa^{-1}=\delta'^{-1}\delta\in D\cap K$,
so this definition is independent of the decomposition.
For multiplication, use
\[
 (\delta\kappa)(\delta'\kappa')
         =\delta\delta'(\delta'^{-1}\kappa\delta')\kappa'.
\]
The conjugation hypothesis makes $\widetilde\chi$ a character.
It extends $\chi$ and is trivial on $D$. The transformation
rules therefore identify $\mathcal A_K^D$ with the corresponding
space $\mathcal A_L$ of level $L$ and character
$\widetilde\chi$.

The projection
$\Gamma\backslash G/K\to\Gamma\backslash G/L$ is onto.
For a representative $x$ of a target class, its fiber consists of
the classes of $x\delta$, $\delta\in D$. Two such classes
coincide exactly when $\delta_1^{-1}\delta_2\in D\cap K$.
In fact, an equality gives
$\gamma x\delta_1\kappa=x\delta_2$ and hence
$x^{-1}\gamma x\in L\subseteq K_0$; equation
\eqref{cd-neat} forces $\gamma=1$.
The converse is immediate. Every fiber thus has cardinality
$[D:D\cap K]$, which is finite because the source is finite and
the target is nonempty. Evaluation on classes, exactly as in
\eqref{cd-column-extension} but with values in $W$, gives
\[
 \dim_k\mathcal A_K=\dim_k W\,|\Gamma\backslash G/K|,
 \qquad
 \dim_k\mathcal A_L=\dim_k W\,|\Gamma\backslash G/L|.
\]
The fiber count proves \eqref{cd-fixed-dimension-eq}.
\end{proof}

\subsection{The coefficient dimension from compact branching}

For a dominant integral tuple $\lambda$ of length $k$, let
$V_\lambda$ be the representation of Section~\ref{cb-section}.
Put $s_k=\prod_{j=0}^{k-1}j!$, with $s_0=1$.

\begin{proposition}[Weyl dimension formula]\label{cd-weyl-dimension}
For $k\ge1$,
\begin{equation}\label{cd-weyl-formula}
        \dim V_\lambda
            =\frac{\prod_{i<j}(\lambda_i-\lambda_j+j-i)}{s_k}.
\end{equation}
Equivalently, for $u_i=\lambda_i+(k+1)/2-i$,
$\dim V_\lambda=\prod_{i<j}(u_i-u_j)/s_k$.
The dimension is the same over any characteristic-zero coefficient
field in which the polynomial representation is defined.
\end{proposition}
\begin{proof}
A common determinant twist changes neither side. Subtract the
smallest entry of $\lambda$ so that all entries are nonnegative.
For $k=1$ the module is one-dimensional and the formula holds.
By the already proved compact branching law,
\[
 \dim V_\lambda
    =\sum_{\lambda_i\ge\mu_i\ge\lambda_{i+1}}\dim V_\mu.
\]
Assume the formula for length $k-1$. Put
$a_i=\lambda_i+k-i$ and $b_i=\mu_i+k-1-i$.
The summation ranges are independent integer intervals
\[
                          a_{i+1}\le b_i<a_i
                                        \qquad(1\le i<k).
\]
The $a_i$ are strictly decreasing nonnegative integers.
For $r=k-1$, the induction hypothesis equals
\[
       \dim V_\mu
            =\det\left[\binom{b_i}{r-j}\right]_{1\le i,j\le r}.
\]
To verify this determinant formula, the polynomial
$\binom{t}{m}$ has degree $m$ and leading coefficient $1/m!$.
Changing the columns from these polynomials to monomials is
triangular. The Vandermonde determinant then gives
$\prod_{i<j}(b_i-b_j)/\prod_{j=0}^{r-1}j!$, with the displayed
sign because the degrees run in decreasing order.

Multilinearity in the rows and the telescoping identity
\[
 \sum_{b=A}^{B-1}\binom bq
                         =\binom B{q+1}-\binom A{q+1}
\]
show that the sum over $\mu$ is
\[
 \det\left[
      \binom{a_i}{r-j+1}-\binom{a_{i+1}}{r-j+1}
              \right]_{1\le i,j\le r}.
\]
This is exactly
$\det[\binom{a_i}{k-j}]_{1\le i,j\le k}$: subtract row
$i+1$ from row $i$ for $i<k$, leaving the last row unchanged.
The row transformation has determinant one; the last column
then has zero entries except for its final entry one.
Applying the same Vandermonde calculation proves
\eqref{cd-weyl-formula}. The assertion about $u$ follows by
subtracting two of its entries. Finally the polynomial modules
are kernels of finite rational linear systems, so extension of a
characteristic-zero coefficient field preserves their dimensions.
\end{proof}

\subsection{Counting the actual decorated lattice classes}

Let $E$ be the CM field of degree $2f$ used in the construction,
and fix a positive definite Hermitian space of dimension $k\ge1$
over $E$. Consider its unimodular $\mathcal O_E$-lattices.
Choose one lattice $N_c$ in each isometry class and put
\begin{equation}\label{cd-mass-definition}
                   M_k=\sum_c\frac1{|\operatorname{Aut}(N_c)|}.
\end{equation}
The class set and the automorphism groups are finite by the
results of Section~\ref{al-section}.

In the actual application, $E/F$ is unramified at every finite
prime and the fixed space is $E^k$ with
$h(x,y)=\sum_i\bar x_i y_i$. Lemma~\ref{acl-local-norm} and
Proposition~\ref{acl-local-basis} show that every unimodular
lattice has an orthonormal basis at every finite completion,
including the split places. All these lattices therefore belong
to the genus of the standard lattice $\mathcal O_E^k$.
Consequently $M_k$ in \eqref{cd-mass-definition} is exactly
the standard genus mass used in Section~\ref{md-section}.

At each of $f$ selected primes of residue cardinality $g_1$ and
each of $f$ selected primes of residue cardinality $g_2$, choose
an isomorphism $N/PN\simeq\mathbb F_{g_i}^{k}$, modulo left
composition by the upper triangular unipotent group
$U_k(\mathbb F_{g_i})$. At each of $f$ further selected primes
of residue cardinality $T$, choose a nonzero linear functional
\[
                         N/PN\longrightarrow\mathbb F_T.
\]
All these primes are distinct. The actual parameters have
$g_1,g_2,T$ prime, with $g_1\ne g_2$ and $T\ne g_i$.
Each selected collection has $f$ elements, and contains only
the chosen primes from the conjugate pairs. A local basis at an
auxiliary level prime is taken modulo the indicated group, not
modulo a Borel subgroup or the scalar matrices.

Let $\mathcal D_T$ be the resulting isometry classes of decorated
lattices, and put
\[
 G_k(q)=\prod_{j=0}^{k-1}(q^k-q^j)
       =q^{k^2}\prod_{j=1}^k(1-q^{-j}),\qquad
 c_k(q)=\frac{G_k(q)}{q^{k(k-1)/2}}.
\]

\begin{lemma}\label{cd-decoration-count}
All decorated stabilizers are trivial, and
\begin{equation}\label{cd-classes-mass}
             |\mathcal D_T|
                 =M_k\bigl(c_k(g_1)c_k(g_2)\bigr)^f(T^k-1)^f.
\end{equation}
\end{lemma}
\begin{proof}
For a full projective rank-$k$ lattice, localization at $P$
is free of rank $k$. Reduction therefore gives
$N/PN\simeq\mathbb F_q^k$, regardless of whether $N$ is
globally free. Its ordered bases form a torsor for
$\operatorname{GL}_k(\mathbb F_q)$. Choosing independent
columns successively gives $G_k(q)$ bases. The group $U_k$
has $q^{k(k-1)/2}$ elements, because its entries strictly above
the diagonal are arbitrary and all other entries are fixed.
Its action on the bases is free, giving exactly $c_k(q)$
choices modulo $U_k$. The dual vector space has $T^k-1$
nonzero vectors, giving that number of nonzero functionals.
Choices at the distinct primes are independent. Thus every
$N_c$ has
\[
                 A=\bigl(c_k(g_1)c_k(g_2)\bigr)^f(T^k-1)^f
\]
possible decorations.

An isometry preserving one decoration has, in the corresponding
local bases, unipotent reductions at both residue characteristics
$g_1,g_2$. It has finite order because it preserves a positive
definite lattice. Proposition~\ref{al-two-primes} makes it the
identity. The same conclusion follows from just the first prime
under the source's additional inequality $2fk<g_1-1$, by the
single-prime argument following that proposition.
The action of $\operatorname{Aut}(N_c)$ on the decorations is
therefore free. There are $A/|\operatorname{Aut}(N_c)|$ orbits
over $N_c$. Summing over $c$ proves \eqref{cd-classes-mass}.
\end{proof}

To incorporate a character of the upper unipotent groups, one
can instead retain the full local bases and let the product of
these groups act along with the rational isometries, with that
character on the coefficient space. The action on the full
decorated objects is free: its rational part is the identity
by the preceding proof, and then the local matrices are the
identity because the bases are isomorphisms. Evaluation therefore
still gives one copy of the coefficient space per class. Thus
the finite-level character does not change the dimension.
The character orientation remains the inverse level convention
proved in Section~\ref{acl-section}; no invariant restriction has
been discarded before proving stabilizer triviality.

\begin{theorem}[The exact source-space dimension]\label{cd-source-dimension}
At each of the $f$ coefficient embeddings let
$\lambda^{(v)}$ be the dominant integral highest weight and
$u_{v,i}=\lambda^{(v)}_i+(k+1)/2-i$. For the actual column
space $\mathcal M_{\epsilon,T}$ with these coefficient modules
and the decorations just specified,
\begin{equation}\label{cd-source-formula}
 \dim\mathcal M_{\epsilon,T}
   =M_k\bigl(c_k(g_1)c_k(g_2)\bigr)^f
           \prod_v\frac{\prod_{i<j}(u_{v,i}-u_{v,j})}{s_k}
                           (T^k-1)^f.
\end{equation}
The same formula holds over the number field of coefficients.
\end{theorem}
\begin{proof}
The class evaluation of Proposition~\ref{acl-class-evaluation},
or equivalently Proposition~\ref{cd-column-decomposition},
gives the number of decorated classes times the dimension of
the coefficient space. That space is the tensor product over
the $f$ embeddings; its dimension is the product of the individual
dimensions. Apply Proposition~\ref{cd-weyl-dimension} and
Lemma~\ref{cd-decoration-count}. Field extension preserves the
dimension, as in Proposition~\ref{cd-weyl-dimension}.
\end{proof}

For the actual weights, the source identifies $u_{v,i}$ with its
Harish--Chandra parameters. The source's zero-based sequence
$\lambda_i-i$ differs from $u_{v,i}$ by one common constant,
which cancels from all differences. Equation
\eqref{cd-source-formula} is therefore precisely
\eqref{as-column-assumption}, with the mass defined by
\eqref{cd-mass-definition}. The integrality of the highest weights
is part of the actual construction; an arbitrary real tuple in
a scalar size estimate need not define a representation.

\subsection{The auxiliary residue-field actions}

Let
$D=\prod_{P\in S_T}(\mathcal O_E/P)^\times$
act by scaling the chosen nonzero functionals, leaving the lattice
and every other decoration unchanged. This action commutes with
isometries and with the changes of local bases.

\begin{lemma}\label{cd-diamond-free}
The action of $D$ on $\mathcal D_T$ is free. After choosing one
class in each $D$-orbit, the column space is, as a $D$-module,
the space of functions $D\to W_0$ with translation action,
for a finite-dimensional coefficient space $W_0$.
\end{lemma}
\begin{proof}
If scaling by $a\in D$ preserves a class, some rational
isometry relates the original decoration to the scaled one.
It preserves the lattice and the local-basis cosets at $g_1,g_2$,
so the preceding stabilizer argument makes it the identity.
Each remaining equality is $a_P\ell_P=\ell_P$ for a nonzero
functional $\ell_P$, forcing $a_P=1$. Thus $a=1$.

Choose representatives of the $D$-orbits and, for every $a\in D$,
use the actual object obtained by scaling such a representative
by $a$. Freeness makes these distinct isometry classes. Evaluation
identifies a column with an arbitrary coefficient vector at each
of them. Scaling permutes these values by translation in $D$;
the lattice and the coefficient vector spaces themselves do not
change. Taking $W_0$ to be the direct sum of the coefficient
spaces at the chosen orbit representatives proves the assertion.
\end{proof}

For a factor $D_P$, let $R_a$ be translation by $a$ on this
function space and put $N_P=\sum_{a\in D_P}R_a$.

\begin{proposition}\label{cd-norm-kernel-dimensions}
The fixed space for one $D_P$ and the simultaneous kernel satisfy
\begin{align}
       \dim\mathcal M_{\epsilon,T}^{D_P}\,(T-1)
                         &=\dim\mathcal M_{\epsilon,T},
                                      \label{cd-one-fixed}\\
       \dim\bigcap_{P\in S_T}\ker N_P\,(T-1)^f
                         &=\dim\mathcal M_{\epsilon,T}\,(T-2)^f.
                                      \label{cd-all-kernels}
\end{align}
The counting argument is valid over every coefficient field.
\end{proposition}
\begin{proof}
On functions on a finite group $D_P$, translation-invariant
functions are exactly the constants. The norm operator sends
a function to the constant function whose value is the sum
of all its values. The sum functional is onto in every
characteristic: a function supported at one point can have
arbitrary sum. Its kernel therefore has dimension $|D_P|-1$.

For several factors, identify the space of functions with
$W_0\otimes_k\bigotimes_P k^{D_P}$. Write
\[
 k^{D_P}=\ker\left(\sum:k^{D_P}\to k\right)\oplus k\delta_1,
\]
where $\delta_1$ is the indicator of the identity element and
has sum one. Each $N_P$ has the same kernel as this sum map
on its tensor factor. The simultaneous kernel is consequently
$W_0$ tensored with the product of all the displayed kernels.
Its dimension is $\dim W_0\prod_P(|D_P|-1)$, whereas the whole
space has dimension $\dim W_0\prod_P|D_P|$.
Finally $|D_P|=T-1$ and $|S_T|=f$. This proves both formulas.
For $T=2$, each factor is trivial and the norm kernel is zero,
exactly as the second formula states.
\end{proof}

The trivial-stabilizer and character conditions in these results
are substantive. For example, take $G=\Gamma=K_0=D$ cyclic of
order two, $K=\{1\}$, and the trivial one-dimensional coefficient
representation. Then $\mathcal A_K$ and its $D$-fixed space both
have dimension one, but $[D:D\cap K]=2$; formula
\eqref{cd-fixed-dimension-eq} would be false without
\eqref{cd-neat}. Separately, take $G=K=K_0=D$ cyclic of order
two, $\Gamma=\{1\}$, and the nontrivial character over
$\mathbb Q$. Now stabilizers are trivial, but $\mathcal A_K$
has dimension one and its $D$-fixed space is zero, so triviality
of $\chi$ on $D\cap K$ cannot be omitted. Both examples are
excluded by the actual hypotheses.

These arguments prove the source-space dimension and the
finite-group fixed-space counts. They do not evaluate the mass
\eqref{cd-mass-definition}, its recurrence in $k$, an
archimedean formal degree, or an ambient automorphic dimension.


